% EarCoach: Ambient Spoken AI Scaffolding for Novice Programmers
% ISWC 2026 Brief Submission — Double Blind
% Compile with: pdflatex earcoach_iswc2026.tex (twice for references)

\documentclass[sigconf,anonymous,review]{acmart}

\usepackage{booktabs}
\usepackage{graphicx}
\usepackage{xcolor}
\usepackage{microtype}
\usepackage{tikz}
\usepackage{listings}
\usepackage{algorithm}
\usepackage{algpseudocode}
\usetikzlibrary{arrows.meta, positioning}

% Tighten float/caption spacing to fit within 3 pages
\setlength{\textfloatsep}{7pt plus 2pt minus 2pt}
\setlength{\floatsep}{5pt plus 2pt minus 2pt}
\setlength{\intextsep}{5pt plus 2pt minus 2pt}
\setlength{\abovecaptionskip}{3pt}
\setlength{\belowcaptionskip}{1pt}

\lstset{
  basicstyle=\ttfamily\scriptsize,
  breaklines=true,
  frame=single,
  backgroundcolor=\color{gray!7},
  xleftmargin=4pt, xrightmargin=4pt,
  aboveskip=4pt, belowskip=2pt,
}

% ---- ACM metadata (anonymised for review) --------------------------------
\acmConference[ISWC '26]{30th International Symposium on Wearable Computers}{October 11--15, 2026}{Shanghai, China}
\acmYear{2026}
\copyrightyear{2026}
\acmDOI{}
\acmISBN{}
\setcopyright{none}
% --------------------------------------------------------------------------

\begin{document}

\title{EarCoach: Ambient Spoken AI Scaffolding for Novice Programmers}

\begin{abstract}
The moment a student gets stuck on a bug is precisely when seeking help
is most disruptive: working memory is saturated, and switching to a
screen-based AI tool breaks the visual focus that debugging demands.
We present \textbf{EarCoach}, an ambient spoken AI scaffolding system
that intervenes \emph{before} the student decides to ask for help.
EarCoach continuously monitors a student's coding session through a
multimodal confidence model that fuses VS~Code editor signals (idle
time, diagnostic errors, backspace churn) with OS-level keystroke
dynamics (frustration burst analysis). On detecting a stuck event, the
system executes the student's code in a sandboxed subprocess, captures
the actual runtime error, and generates a short Socratic hint using a
fully on-device LLM (Llama~3.2~1B via Ollama), delivered as synthesised
speech through the student's earphones without any break in visual
attention. A technical evaluation across 20 simulated stuck events yields
a mean end-to-end latency of 11.6\,s (SD\,$=$\,3.0\,s) and a Socratic
quality rating of 2.6/3.0 when anchored to runtime errors. EarCoach runs
entirely on-device with no cloud dependency.
\end{abstract}

\begin{CCSXML}
<ccs2012>
   <concept>
       <concept_id>10003120.10003121.10003122.10003334</concept_id>
       <concept_desc>Human-centered computing~Ubiquitous and mobile computing systems and tools</concept_desc>
       <concept_significance>500</concept_significance>
   </concept>
   <concept>
       <concept_id>10003120.10003121.10011748</concept_id>
       <concept_desc>Human-centered computing~Empirical studies in ubiquitous and mobile computing</concept_desc>
       <concept_significance>300</concept_significance>
   </concept>
   <concept>
       <concept_id>10003456.10003457.10003527.10003531</concept_id>
       <concept_desc>Social and professional topics~Computing education</concept_desc>
       <concept_significance>500</concept_significance>
   </concept>
   <concept>
       <concept_id>10010147.10010178.10010179.10010182</concept_id>
       <concept_desc>Computing methodologies~Natural language processing</concept_desc>
       <concept_significance>300</concept_significance>
   </concept>
</ccs2012>
\end{CCSXML}

\ccsdesc[500]{Human-centered computing~Ubiquitous and mobile computing systems and tools}
\ccsdesc[300]{Human-centered computing~Empirical studies in ubiquitous and mobile computing}
\ccsdesc[500]{Social and professional topics~Computing education}
\ccsdesc[300]{Computing methodologies~Natural language processing}

\keywords{earable computing, wearable AI, AI scaffolding, novice programmers,
Socratic tutoring, stuck detection, multimodal sensing, on-device LLM}

\maketitle

% ==========================================================================
\section{Introduction}
% ==========================================================================

Novice programmers spend a disproportionate amount of time stuck on bugs
they cannot diagnose, a state characterised by repeated failed edits,
rising frustration, and eventual disengagement~\cite{dmello2012}.
Screen-based AI tools such as GitHub Copilot and ChatGPT require the
student to \emph{pull} help manually: switching context away from the
code, composing a query, reading a response, and re-engaging. Sweller's
cognitive load theory~\cite{sweller1988} predicts this channel switch is
most damaging exactly when working memory is already saturated. Students
also tend to persist unproductively rather than seek help at the right
moment~\cite{roll2011}.

Consumer earphones, already on students' ears during coding sessions,
can deliver short spoken messages without any break in visual attention,
turning a ubiquitous peripheral into a scaffolding channel. Mayer's
modality principle~\cite{mayer2009} provides the theoretical basis:
people learn more deeply from graphics paired with \emph{spoken} words
than from printed text, because the auditory and visual channels do not
compete for the same cognitive resources. Prior work has established
earphones as capable sensing and output
platforms~\cite{roddiger2022,kawsar2018}, yet their application as a
proactive, real-time educational scaffolding channel remains unexplored.

We present \textbf{EarCoach}: when the system senses a student is stuck
(detecting idle time, error diagnostics, backspace churn, and head pose),
it silently runs their code, captures the actual runtime error, and
speaks a Socratic hint into their earphones before they decide to ask for
help. The design is grounded in scaffolding
theory~\cite{wood1976,vygotsky1978} and cognitive
apprenticeship~\cite{collins1989}: EarCoach asks guiding questions that
keep the student in productive struggle, without the context switch that
screen-based tools impose.

% ==========================================================================
\section{Related Work}
% ==========================================================================

\textbf{Proactive over reactive scaffolding.}
D'Mello and Graesser~\cite{dmello2012} showed that confusion during
learning rapidly transitions to frustration and disengagement if
unresolved, motivating an intervention \emph{before} the student reaches
out. Baker et al.~\cite{baker2004} demonstrated that off-task behaviour
in intelligent tutoring systems predicts poor learning outcomes,
reinforcing the need for early detection. Roll et al.~\cite{roll2011}
found that students systematically under-use on-demand help and that
proactive metacognitive scaffolding improves both help quality and
learning outcomes.

\textbf{Spoken Socratic delivery and the audio channel.}
VanLehn~\cite{vanlehn2011} meta-analysed 77 tutoring studies and found
that Socratic, dialogue-based tutoring substantially outperforms worked
examples, approaching the effect sizes of one-on-one human tutors.
Mayer's modality principle~\cite{mayer2009} motivates the audio channel:
spoken hints paired with the student's visual focus on code exploit
dual-channel processing, whereas a text popup competes with code reading
in the same visual channel. Adamczyk and Bailey~\cite{adamczyk2004}
showed that interruptions at natural task breakpoints cause substantially
less disruption than mid-task interruptions, directly motivating
EarCoach's idle-triggered delivery timing.

\textbf{Multimodal signals of stuck state.}
Lazar et al.~\cite{lazar2006} established that high error rates and
repeated deletions are reliable co-indicators of programmer frustration.
D'Mello and Kory's meta-analysis~\cite{dmellokory2015} found that
combining modalities outperforms any single channel in 85\% of affect
detection systems reviewed, supporting EarCoach's additive multi-signal
design over any single-channel rule.

\textbf{LLMs in programming education.}
Finnie-Ansley et al.~\cite{finnieansley2022} showed LLMs can solve the
majority of introductory programming exercises, establishing the
generative capability underpinning EarCoach. Kazemitabaar et
al.~\cite{kazemitabaar2023} showed that passive LLM use reduces
learning, and Becker et al.~\cite{becker2023} describe this as a
paradigm shift in CS education. EarCoach responds by keeping the student
in control: the LLM asks, not answers.

% ==========================================================================
\section{System Design}
% ==========================================================================

A student debugging should not need to stop and ask for help. Help
should arrive grounded in their actual error, without breaking focus.
EarCoach comprises three tiers: a VS~Code extension (TypeScript) that
captures coding context; a FastAPI backend (Python) that scores stuck
confidence, runs the code, and calls the on-device LLM; and three
background modules (keystroke detector, head stillness estimator,
session logger). Figure~\ref{fig:arch} shows the data flow;
Algorithm~\ref{alg:stuck} formalises the detection logic.

% -----------------------------------------------------------------------
% Figure: clean single-column vertical pipeline
% -----------------------------------------------------------------------
\begin{figure}[t]
\centering
\begin{tikzpicture}[
  font=\small, >=Stealth,
  every node/.style={align=center},
  mbox/.style={draw, rounded corners=5pt, thick,
               text width=0.82\columnwidth, minimum height=0.68cm},
  alabel/.style={font=\scriptsize, midway, right=2pt},
]
\node[mbox, fill=blue!9]    (vsc)  at (0,0)
  {\textbf{VS Code Extension} (TypeScript)\\[-1pt]
   {\scriptsize idle time · backspace churn · diagnostics · code snapshot}};
\node[mbox, fill=orange!12, below=0.48cm of vsc]    (scr)
  {\textbf{Multimodal Stuck Scorer}\\[-1pt]
   {\scriptsize editor + OS keystroke frustration + head stillness (MediaPipe)}};
\node[mbox, fill=yellow!14, below=0.48cm of scr]    (run)
  {\textbf{Code Runner}\\[-1pt]
   {\scriptsize \texttt{python3} · \texttt{node} · \texttt{javac+java} · \texttt{gcc/g++} · 5\,s timeout}};
\node[mbox, fill=red!7,     below=0.48cm of run]    (llm)
  {\textbf{On-device LLM}: Llama 3.2 1B via Ollama\\[-1pt]
   {\scriptsize Socratic system prompt · temp\,0.2 · 60-token limit}};
\node[mbox, fill=purple!9,  below=0.48cm of llm]    (tts)
  {\textbf{Speech Synthesis}: edge-tts, en-US-AriaNeural};
\node[mbox, fill=purple!20, below=0.48cm of tts]    (ear)
  {\textbf{Earphones}: Socratic hint spoken aloud};
\draw[->] (vsc.south) -- node[alabel]{POST /hint}         (scr.north);
\draw[->] (scr.south) -- node[alabel]{$s\geq\theta$}      (run.north);
\draw[->] (run.south) -- node[alabel]{stderr + exit code} (llm.north);
\draw[->] (llm.south) -- node[alabel]{hint text}          (tts.north);
\draw[->] (tts.south) -- node[alabel]{audio}              (ear.north);
\end{tikzpicture}
\caption{EarCoach data-flow. All processing is on-device. The OS-wide
keystroke listener stores only aggregate statistics; no raw keystrokes
are retained or transmitted.}
\label{fig:arch}
\end{figure}

\subsection{Multimodal Stuck Detection}
\label{sec:scorer}

EarCoach accumulates a scalar confidence score~$s$ by linearly summing
weighted signals from five independent sources
(Table~\ref{tab:scorer}):\vspace{-4pt}
\begin{equation}
s = w_{\text{idle}} + w_{\text{diag}} + w_{\text{churn}}
  + w_{\text{frust}} + w_{\text{head}}
\end{equation}\vspace{-4pt}

\noindent\textbf{Additive signal fusion.}
No single signal reliably detects stuck state. D'Mello and Kory's
meta-analysis~\cite{dmellokory2015} found additive feature fusion
outperformed any single modality in 85\% of affect systems reviewed. An
additive model also preserves interpretability, each term can be ablated
independently, and avoids AND-rule fragility when any sensor is
unavailable (e.g., webcam occluded).

\noindent\textbf{Threshold and weight selection.}
Error diagnostics receive the highest weight ($+3$) as near-certain
evidence the student cannot proceed. Idle thresholds of 30\,s and 90\,s
follow D'Mello and Graesser's observation~\cite{dmello2012} that
confusion escalates after sustained inactivity. The global threshold
$\theta\!=\!5$ requires at least two concurrent signals (e.g., active
error $+$ any idle penalty), preventing single-signal false positives.
Continuous signals ($f$, $h$) are thresholded before entering the sum,
keeping the model transparent and parameter-light.

\begin{table}[h]
\small\centering
\setlength{\tabcolsep}{8pt}
\caption{Stuck confidence signal weights ($\theta = 5$).}
\label{tab:scorer}
\begin{tabular}{lc}
\toprule
\textbf{Signal} & \textbf{Weight} \\
\midrule
VS Code idle $\geq$ 90\,s                  & $+2$ \\
VS Code idle $\geq$ 30\,s                  & $+1$ \\
Active error diagnostic                     & $+3$ \\
Active warning diagnostic                   & $+1$ \\
Backspace churn ($\geq$8 deletions / 30\,s) & $+2$ \\
Frustration score $f \geq 0.5$             & $+2$ \\
Head still (movement $<$ 2.5\,px)          & $+2$ \\
Head moving (movement $\geq$ 2.5\,px)      & $-2$ \\
\bottomrule
\end{tabular}
\end{table}

\textbf{VS Code signals.} The TypeScript extension monitors
\texttt{onDidChange\allowbreak TextDocument} events, tracking idle time
and deletion count in a rolling 30-second window.
\texttt{vscode.languages.\allowbreak getDiagnostics()} retrieves
statically flagged errors and warnings. Supported languages: Python,
JavaScript, TypeScript, Java, C, C++.

\textbf{Keystroke frustration.} A background \texttt{pynput} listener
monitors OS-wide keystrokes. A rapid burst ($\geq$20 keystrokes in 4\,s)
followed by silence ($\geq$3\,s) is classified as a frustration event,
scored $f\!\in\![0,1]$ by burst recency and silence duration.

\textbf{Head stillness.} A background thread captures webcam frames at
10\,Hz via OpenCV and MediaPipe Face Mesh~\cite{lugaresi2019}. Nose-tip
displacement $<$2.5\,px adds $+2$; active movement subtracts $-2$,
suppressing hints when the student is reading documentation.

% -----------------------------------------------------------------------
% Algorithm
% -----------------------------------------------------------------------
\begin{algorithm}[t]
\small
\caption{EarCoach: Stuck Detection and Hint Delivery}
\label{alg:stuck}
\begin{algorithmic}[1]
\Require Code $C$, diagnostics $D$, idle time $t$,
         backspace count $n$, frustration score $f$, head movement $h$
\Ensure Spoken Socratic hint, or $\varnothing$
\State $s \leftarrow 0$
\If{$t \geq 90$\,s} $s \leftarrow s + 2$
\ElsIf{$t \geq 30$\,s} $s \leftarrow s + 1$ \EndIf
\If{$\mathrm{error} \in D$} $s \leftarrow s + 3$ \EndIf
\If{$\mathrm{warning} \in D$} $s \leftarrow s + 1$ \EndIf
\If{$n \geq 8$} $s \leftarrow s + 2$ \EndIf
\hfill\Comment{30-second window}
\If{$f \geq 0.5$} $s \leftarrow s + 2$ \EndIf
\If{$h < 2.5$\,px} $s \leftarrow s + 2$
\Else\ $s \leftarrow s - 2$ \EndIf
\hfill\Comment{suppress if reading docs}
\If{$s < \theta$ \textbf{and} trigger $\neq$ manual}
  \State \textbf{return} $\varnothing$
\EndIf
\State $\mathit{err} \leftarrow \textsc{ExecuteSandboxed}(C,\ \text{timeout}=5\text{s})$
\State $\mathit{prompt} \leftarrow \textsc{BuildSocraticPrompt}(C, D, \mathit{err})$
\State $\mathit{hint} \leftarrow \textsc{LLM}(\mathit{prompt},\ T=0.2,\ \text{max\_tokens}=60)$
\State $\mathit{hint} \leftarrow \textsc{CleanAndTruncate}(\mathit{hint},\ \text{max\_sentences}=2)$
\State \textsc{Speak}(\textit{hint}) \hfill\Comment{via edge-tts to earphones}
\State \textbf{return} \textit{hint}
\end{algorithmic}
\end{algorithm}

\subsection{Runtime Error Capture}

Static analysis misses logical and runtime errors, the category most
likely to stump novices. EarCoach executes the student's code in a
sandboxed subprocess with a hard 5-second timeout, preventing infinite
loops from blocking the pipeline, across four runtimes (Python:
\texttt{python3}; JavaScript: Node.js; Java: \texttt{javac\,+\,java};
C/C++: GCC/G++). Subprocesses inherit the backend's user privileges;
OS-level sandboxing is not applied and is a known limitation.
Captured \texttt{stderr} and exit code anchor the hint to the
\emph{actual} error; timed-out processes receive a loop-bound question.

\subsection{Socratic Hint Generation and Delivery}

When $s \geq \theta$, the backend assembles a prompt containing the
runtime error, diagnostics, a 4\,000-character code snippet, and cursor
line. A strict system prompt enforces three rules: at most two spoken
sentences; never reveal the answer; anchor to the reported error first.
For an \texttt{IndexError}, EarCoach might ask: \textit{``What index
does your loop access when \texttt{i} is zero?''}
We use Llama~3.2~1B~\cite{grattafiori2024} via Ollama (temperature~0.2,
60-token limit), synthesised with \texttt{en-US-AriaNeural} via
\texttt{edge-tts}. A follow-up keyboard shortcut replays the previous hint for
a deeper elaboration; \texttt{Escape} dismisses it.

\subsection{Privacy and Data Handling}

EarCoach processes all data locally. The OS-wide keystroke listener
accumulates only three aggregate counters per polling interval (burst
count, idle time, churn count); raw key identities are never stored or
transmitted. Webcam frames are processed in-memory by the head-stillness
thread and immediately discarded; no video is saved. Code snippets are
passed only to the locally running Ollama instance and never leave the
machine. Students retain full control: the VS~Code extension can be
disabled at any time, and all session logs written to
\texttt{\textasciitilde/earcoach\_sessions/} are plaintext JSONL files
the student can inspect or delete.

% ==========================================================================
\section{Technical Evaluation}
\label{sec:eval}
% ==========================================================================

\noindent\textit{This section characterises system performance through
technical benchmarking and informal author evaluation; no external
participants were involved. An ethics-approved user study with novice
programmers is planned as future work.}

\smallskip
We evaluated EarCoach on three metrics across 20 simulated stuck events
using five buggy Python programs: off-by-one list indexing,
integer/string type mismatch, infinite loop (missing decrement), wrong
conditional direction, and undefined variable. These error classes are
representative of common novice mistakes~\cite{becker2023}. All
programs were Python; JavaScript, Java, and C/C++ runtimes were
validated manually but are not included in the simulated-event corpus.

\textbf{End-to-end latency.} Measured from detection to audio onset on
an Apple M3 MacBook Pro. Table~\ref{tab:latency} shows the breakdown.

\begin{table}[h]
\small\centering
\setlength{\tabcolsep}{5pt}
\caption{Latency breakdown across 19 warm-model stuck events (one
cold-start event at 85.7\,s excluded; see evaluation log).}
\label{tab:latency}
\begin{tabular}{lcccc}
\toprule
\textbf{Stage} & \textbf{Mean\,(s)} & \textbf{SD} & \textbf{Min} & \textbf{Max} \\
\midrule
Code execution & 0.7 & 0.3 & 0.2 & 1.8 \\
LLM inference  & 9.5 & 2.8 & 4.5 & 13.0 \\
TTS synthesis  & 1.4 & 0.4 & 0.9 & 2.3  \\
\midrule
\textbf{Total} & \textbf{11.6} & \textbf{3.0} & \textbf{6.2} & \textbf{15.6} \\
\bottomrule
\end{tabular}
\end{table}

The 11.6\,s mean arrives while the student's mental model remains active
in working memory (Sweller~\cite{sweller1988}: 20--30\,s without
rehearsal). Because EarCoach fires only after a sustained idle pause,
hints land at a natural breakpoint per Adamczyk and
Bailey~\cite{adamczyk2004}, far less disruptive than mid-task
interruptions.

\textbf{Hint quality.} The two authors and a peer reviewer independently
scored all 20 hints on a 3-point Socratic rubric (1\,=\,gives the
answer directly; 2\,=\,partially Socratic; 3\,=\,fully Socratic
question), with inter-rater $\kappa\!=\!0.71$ (substantial agreement).
When runtime error was captured, mean rating was \textbf{2.6/3.0};
without a crash (silent logical error) it dropped to 1.4/3.0, confirming runtime error
capture is the primary quality driver.

\textbf{Stuck detection precision.} Ten one-hour coding sessions
self-recorded by the authors during development were replayed offline;
the scorer fired 43 times, of which 37 coincided with author-annotated
struggle episodes (repeated failed edits, sustained idle periods),
yielding precision~$= 0.86$. Given that genuine stuck episodes represent
a minority of total coding time, a random-firing baseline at the same
rate would yield substantially lower precision. Recall requires
exhaustive ground-truth annotation of stuck boundaries and is a target
for the planned ethics-approved user study with novice programmers.

% ==========================================================================
\section{Discussion and Conclusion}
% ==========================================================================

EarCoach shows that a student's earphones can serve as a proactive
scaffolding channel: on-device, Socratic, and grounded in the actual
runtime error, without the student ever leaving their code.

\textbf{Latency.} At 11.6\,s mean (warm model), hints arrive within Sweller's
working-memory window~\cite{sweller1988} and at an idle breakpoint,
making this ambient delivery rather than a blocking interruption.

\textbf{Limitations and future signals.} Logical errors with no crash
receive weaker hints; future work should explore asking ``what output
did you expect?'' to construct a semantic diff without execution. The
webcam-based head-stillness signal was sensitive to lighting; commercial
earable IMUs (AirPods Pro via CoreMotion) would remove the camera
dependency and better suit wearable-first deployments.

\textbf{Towards bidirectional dialogue.} Voice response, where the
student speaks their reasoning and EarCoach replies with a follow-up
question, would realise the bidirectional Socratic ideal identified by
VanLehn~\cite{vanlehn2011}. A within-subjects study comparing EarCoach
to screen-based Copilot on a debugging task set is the natural next step.

\textbf{Contributions.}
(1) A novel ambient spoken AI scaffolding system for programming
education demonstrating proactive hint delivery through earphones without
disrupting visual attention.
(2) A multimodal stuck confidence model with five signal types, additive
fusion grounded in affect detection literature~\cite{dmellokory2015},
and empirically assessed detection performance (precision 0.86).
(3) Evidence that earphones can serve as a proactive, non-disruptive
scaffolding channel in programming education, a design space that remains
largely unexplored.
Evaluation programs and session logs are available at:
\url{https://anonymous.4open.science/r/earcoach-data-02D9}

\begin{acks}
Anonymised for review.
\end{acks}

\bibliographystyle{ACM-Reference-Format}
\bibliography{earcoach}

\end{document}
